\subsection{切空间、切线性映射}

5.5.1. 设 X 是一个流形，且 \(a\in\mathbb{X}\)。考虑有序对 \((c,h)\)，其中 \(c=(\mathrm{U},\varphi,\mathrm{E})\) 是流形 X 在 \(a\) 处的一个图，\(h\) 是 E 的一个元素。若映射 \((c,h)\) 在 \((c',h')\) 处的导数（它定义在 \(\varphi(a)\) 的某个邻域上）将 \(\varphi'\circ\varphi^{-1}\) 变为 \(\varphi(a)\)，则称两个这样的有序对 \(h\) 和 \(h'\) 等价。这样就在有序对 \((c,h)\) 之间得到一个等价关系，并把 X 在 \(a\) 处的切向量定义为等价有序对 \((c,h)\) 的等价类（集合论，第~II~章，§ 6，第 9 号）。

X 在 \(a\) 处的切向量构成一个集合，记为 \(\mathrm{T}_{a}(\mathrm{X})\)。若 \(c = (\mathrm{U}, \varphi, \mathrm{E})\) 是流形 X 在 \(a\) 处的一个图，则把 E 中元素 \(\theta_{c}\) 与由有序对 \(\mathrm{T}_{a}(\mathrm{X})\) 表示的切向量对应起来的映射 \(h\) 从 E 到 \((c, h)\) 是双射。借助 \(\mathrm{T}_{a}(\mathrm{X})\)，把 E 上的拓扑 K-向量空间结构传递给 \(\theta_{c}\)，所得结构与 \(c\) 的选择无关，并使 \(\mathrm{T}_{a}(\mathrm{X})\) 成为一个 Banach 空间，称为 X 在 \(a\) 处的切空间。\(+\infty\) 的维数（有限或 \(\mathrm{T}_{a}(\mathrm{X})\)）等于 X 在 \(a\) 处的维数。

5.5.2. 设 X、Y 是两个流形，\(f\) 是从 X 到 Y 的态射，\(a\) 是 X 的一点。考虑 X 在 a 处的图 \(c = (U, \varphi, E)\)，以及 Y 在 \(a\) 处的图 \(c' = (V, \psi, F)\)，其中 \(b\)\(f(U) \subset V\)；从 \(\Phi = \psi \circ f \circ \varphi^{-1}\)\(\varphi(U)\) 到 \(F\) 的映射 \(C^{r}\) 是 C＾r 类的，并且其在点  处的导数 \(\mathbf{D}\Phi(\varphi(a))\)\(\varphi(a)\)

是从 E 到 F 的连续线性映射。从 \(\theta_{c'} \circ D\Phi(\varphi(a)) \circ \theta_{c}^{-1}\) 到 \(T_{a}(X)\) 的连续线性映射 \(T_{b}(Y)\) 不依赖于所选的图 \(c\) 和 \(c'\)；记作 \(T_{a}(f)\)，称为 \(f\) 在 \(a\) 处的切线性映射。若 \(g\) 是从 Y 到流形 Z 的态射，则有：

\[
\mathrm{T} _ {a} (g \circ f) = \mathrm{T} _ {f (a)} (g) \circ \mathrm{T} _ {a} (f).
\]

5.5.3. 设 \(f:X\to Y\) 是流形的一个态射。若 \(f\) 局部为常值，则对 X 中所有 \(\mathrm{T}_{a}(f)=0\)，都有 \(a\)。反之，若对所有 \(\mathrm{T}_{a}(f)=0\) 都有 \(a\in X\)，且域 K 的特征为 0，则 \(f\) 局部为常值。

5.5.4. 设 U 是流形 X 的一个开子集，f 是从 U 到 X 的典范嵌入；对于 U 中每个点 a，从 \(\mathrm{T}_{a}(f)\) 到 \(\mathrm{T}_{a}(\mathrm{U})\) 的映射 \(\mathrm{T}_{a}(\mathrm{X})\) 是拓扑向量空间的同构，今后据此把这两个空间等同起来。

5.5.5. 设 U 是 Banach 空间 E 的一个开子集；若 \(\iota\) 是 U 到 E 的典范嵌入，则三元组 \(c = (U, \iota, E)\) 是 U 的一个图，而图册 \(\{c\}\) 定义了 U 的流形结构。给定 U 中一点 \(a\)，映射 \(\theta_c\) 是从 E 到 \(T_a(U)\) 的拓扑向量空间同构；其逆同构记为 \(\zeta_E(a)\)。

设 \(f\) 是从 U 到 Banach 空间 F 的开子集 V 的 \(\mathbf{C}^{r}\) 类映射；对于 U 中每一点 \(a\)，有：

\[
\mathbf {D} f (a) \circ \zeta_ {\mathbf {E}} (a) = \zeta_ {\mathbf {F}} (f (a)) \circ \mathrm{T} _ {a} (f)
\]

其中 Df(a) 是 f 在 a 处的导数（1.2.1）。换言之，下图：

\[
\begin{array}{c} \mathrm{T} _ {a} (\mathrm{U}) \xrightarrow {\mathrm{T} _ {a} (f)} \mathrm{T} _ {f (a)} (\mathrm{V}) \\ \Biggl \downarrow \zeta_ {\mathbf {E}} (a) \qquad \qquad \qquad \Biggl \downarrow \zeta_ {\mathbf {F}} (f (a)) \\ \mathrm{E} \xrightarrow {\mathrm{D} f (a)} \mathrm{F} \end{array}
\]

是交换的。

5.5.6. X 在 a 处的余切向量，也称为切余向量，或简称为 a 处的余向量，是拓扑向量空间 \(\mathrm{T}_{a}(\mathrm{X})\) 上的任意连续线性形式；因此，这些余向量就是 X 在 a 处的切空间的拓扑对偶 \(\mathrm{T}_{a}(\mathrm{X})'\) 的元素。赋予这个空间以 \(\mathrm{T}_{a}(\mathrm{X})\) 的有界子集上一致收敛的拓扑；在此拓扑下，\(\mathrm{T}_{a}(\mathrm{X})'\) 是一个 Banach 空间，称为 X 在 a 处的余切空间。它也记为 \(\mathrm{T}_{a}^{\prime}(\mathrm{X})\)。

设 \(f\) 是从 X 到 Banach 空间 E 的态射。称从 \(f\) 到 E 的连续线性映射 \(a\) 为 \(d_{a}f\) 在 \(df_{a}\) 处的微分，记作 \(\zeta_{\mathbf{E}}(f(a)) \circ \mathbf{T}_{a}(f)\) 或 \(\mathbf{T}_{a}(\mathbf{X})\)。对于 \(t \in \mathbf{T}_{a}(\mathbf{X})\)，有时把线性映射 \(t(f)\) 在点 t 处的值 \(d_{a}f(t)\) 记为 \(d_{a}f\) 或 t.f。映射 \(f \mapsto d_{a}f\) 是线性的。

若 \(E = K\)，则 \(d_a f\)\(f\) 在 \(a\) 处的微分 \(X\)\(a\) 是 X 在 a 处的一个余向量。

设 E、F、G 是三个 Banach 空间，且 \((u, v) \mapsto u \cdot v\) 是从 \(F \times G\) 到 E 的连续双线性映射。设 f（分别为 g）是从 X 到 F（分别为 G）的态射。则映射 \(f \cdot g : x \mapsto f(x) \cdot g(x)\) 是从 X 到 E 的态射，并且对于每个 \(t \in \mathrm{T}_{a}(\mathrm{X})\)，有：

\[
d _ {a} (f. g) (t) = f (a). d _ {a} g (t) + d _ {a} f (t). g (a).
\]

特别取 G = E、F = K，且所讨论的双线性映射为乘法，则有：

\[
d _ {a} (f g) = f (a) d _ {a} g + g (a) d _ {a} f.
\]

5.5.7. 设 X 是一个局部有限维流形，\(\xi^{1},\ldots,\xi^{m}\) 是定义在 X 中一点 \(a\) 的某个开邻域 U 上的态射函数。下列条件等价：

\begin{enumerate}
\def\labelenumi{(\roman{enumi})}
\item
  存在包含于 U 的 a 的一个开邻域 V，使得函数 \(\xi^{i}|V\)（对 \(1 \leqslant i \leqslant m\)）构成 X 在 V 中的一个坐标系；
\item
  微分 \(d_a\xi^i\)（对 \(1 \leqslant i \leqslant m\)）构成 \(T_a(X)'\) 的一个基；
\item
  映射 \(\xi = (\xi^1, \ldots, \xi^m)\) 从 U 到 \(K^m\) 在 \(a\) 处是 étale 的（参见第 5.7.6 号）。
\end{enumerate}

微分 \(d_{a}\xi^{1},\ldots,d_{a}\xi^{m}\) 线性无关的充分必要条件是：存在包含于 U 的 \(a\) 的一个邻域 V，使得函数 \(\xi^{i}|V\) 是 X 在 V 中某个坐标系的一部分。

5.5.8. 设 X 是一个局部有限维流形，且 \(\xi = (\xi^1, \ldots, \xi^m)\) 是 X 在开集 U 中的一个坐标系。令 \(a \in \mathbf{U}\)：我们把 \((\partial_{1,a}, \ldots, \partial_{m,a})\) 记为切空间 \(\mathrm{T}_a(\mathrm{X})\) 的基，它与 \((d_a\xi^1, \ldots, d_a\xi^m)\) 的基 \(\mathrm{T}_a(\mathrm{X})'\) 对偶。因此有：

\[
\partial_ {i, a} (\xi^ {j}) = \delta_ {i, j} \quad (\text { Kronecker index }).
\]

切向量 \(\partial_{i,a}\) 也记为 \((\partial/\partial\xi^{i})_{a}\)。

设 \(f\) 是 U 上取值于 Banach 空间 E 的 \(\mathbf{C}^r\) 类函数。把函数 \(\partial_i f\) 记为 \(\partial f / \partial \xi^i\) 或 \(a \mapsto \partial_{i,a}(f)\)：它是 U 上取值于 E 的 \(\mathbf{C}^{r-1}\) 类函数（当 \(r = 1\) 时是连续函数）。它在 U 中一点 \(a\) 处的值有时记为 \((\partial f / \partial \xi^i)_a\)。对于每个 \(t \in \mathrm{T}_a(X)\)，有：

\[
d _ {a} f (t) = \sum_ {i = 1} ^ {m} d _ {a} \xi^ {i} (t) \frac {\partial f}{\partial \xi^ {i}} (a),
\]

也可写成：

\[
d _ {a} f = \sum_ {i = 1} ^ {m} \frac {\partial f}{\partial \xi^ {i}} (a) d _ {a} \xi^ {i}.
\]

这些记号与 1.6.3 中的记号一致。

5.5.9. 设 X 和 Y 是两个流形，f 是从 X 到 Y 的态射。称线性映射 \(rg_{a}f\) 的秩（有限或等于 +∞）为 f 在 X 中一点 a 处的秩，记为 \(\overline{T}_{a}(f)\)。映射 \(a \mapsto rg_{a}f\) 是下半连续的。

